\section*{Summary}
\vspace{-2ex}\titlerule\vspace{4ex}

Accurate temperature measurement is critical in industrial metal forming and heat-treatment processes, where pyrometers provide non-contact sensing at extreme temperatures. However, raw pyrometer signals suffer from two key problems that are addressed in this individual thesis contribution: measurement noise introduced by electronic circuitry and infrared reflections from surrounding equipment, and the large volume of raw data generated during continuous production runs that must be stored and transmitted efficiently.

This thesis addresses these two problems through the design, implementation, and evaluation of automated signal denoising and data compression stages within a modular pre-processing pipeline, structured around two acceptance test procedures. In the denoising stage~(ATP-1), three classical methods were implemented and compared: moving average filter~(k\,=\,7), median filter~(k\,=\,7), and Savitzky--Golay filter~(window\,=\,11, degree\,=\,3). These were evaluated against two machine learning approaches: a one-dimensional Convolutional Neural Network~(CNN) with four Conv1d layers and a Long Short-Term Memory~(LSTM) network with two layers and hidden size~64. Testing on the NIST open-source additive manufacturing dataset~(Layer~01) showed that the CNN achieved the best denoising result with a noise standard deviation of~64\,\textdegree C compared to~113\,\textdegree C in the raw signal, representing a noise reduction of~43.4\,\%. Among classical methods, the moving average achieved the highest noise reduction~(39.8\,\%) while the median filter best preserved sharp temperature peaks.

In the compression stage~(ATP-3), three classical methods were evaluated: downsampling~(2$\times$), Truncated SVD~(rank\,=\,10), and Wavelet compression~(Haar transform, 10\,\% coefficient retention). Two machine learning compression methods were additionally implemented: PCA compression~(3~components, window\,=\,32) and Autoencoder compression~(bottleneck\,=\,3, window\,=\,32). Wavelet compression achieved the best balance between compression ratio~(5.0$\times$) and reconstruction accuracy~(RMSE\,=\,36.55\,\textdegree C), satisfying the ATP-3 target of~$CR > 4\times$. PCA and Autoencoder achieved a higher ratio of~10.7$\times$ at greater reconstruction error.

The individual contribution of this thesis is the systematic investigation of signal denoising and data compression methods for pyrometer time-series, demonstrating that the CNN denoiser and Wavelet compression provide the best performance for their respective stages. The pipeline is delivered as a modular, open-source Python implementation suitable for integration into industrial monitoring systems.
